\documentclass[10pt]{article}
\ifdefined\pdfinfoomitdate\pdfinfoomitdate=1\fi
\ifdefined\pdfsuppressptexinfo\pdfsuppressptexinfo=15\fi
\ifdefined\pdftrailerid\pdftrailerid{<A64376B12B724B17A88322FE379B49E4><A64376B12B724B17A88322FE379B49E4>}\fi
\usepackage[T1]{fontenc}
\usepackage[utf8]{inputenc}
\usepackage{lmodern}
\usepackage{amsmath,amssymb,array,longtable,booktabs}
\usepackage[a4paper,margin=20mm]{geometry}
\usepackage[hidelinks]{hyperref}
\hypersetup{pageanchor=false,pdftitle={Apparatus to Paul Gordan, De linea geodetica (1862)},pdfauthor={Paul Gordan; critical apparatus}}
\setlength{\parindent}{0pt}
\setlength{\parskip}{0.40em}
\setlength{\emergencystretch}{2em}
\sloppy
\providecommand{\GordanProjectName}{Paul Gordan, De linea geodetica (1862)}
\newcommand{\gq}[1]{\glqq#1\grqq}
\newcommand{\tg}{\operatorname{tg}}
\newcommand{\arc}{\operatorname{arc}}
\newcommand{\dd}{\,d}
\newcommand{\GordanDitto}{\mbox{\normalfont\ttfamily\char34\kern0.45em\char34}}
\newcommand{\caput}[1]{\par\vspace{0.35em}\begin{center}{\large\bfseries CAPUT #1.}\end{center}\vspace{-0.75em}}
\newcommand{\sect}[1]{\par\vspace{0.28em}\begin{center}{\bfseries\S\ #1.}\end{center}\vspace{-0.65em}}
\newcommand{\GordanUnitStart}[2]{%
  \clearpage\thispagestyle{empty}%
  \phantomsection\hypertarget{PU#1}{}%
  \pdfbookmark[1]{Physical unit #1 -- #2}{bookmark-PU#1}%
}
\newcommand{\GordanTopFolio}[1]{%
  \begin{center}\textbf{#1}\\[-0.45em]\rule{3.0cm}{0.35pt}\end{center}\vspace{0.1em}%
}
\newcommand{\GordanBottomFolio}[1]{\par\vfill\hfill\textbf{#1}\par}
\newcommand{\GordanOrnament}{\begin{center}\rule{2.2cm}{0.35pt}\kern-1.25cm\raisebox{-0.12ex}{\scriptsize$\blacklozenge$}\end{center}}

\begin{document}
\begin{center}
{\Large\bfseries Apparatus to Paul Gordan, \emph{De linea geodetica} (1862)}\\[0.45em]
{\large Integrated candidate: physical units 1--47}
\end{center}

This restrained apparatus is separate from the standalone Latin and English editions. It records one explicit disposition for every physical unit, copy-matter exclusions, and source-critical readings established by direct authority-pixel replay. The diplomatic Latin preserves printed anomalies; the faithful English retains their mathematical content without silently substituting conjectural corrections.

\section*{Unit dispositions}
\small
\begin{longtable}{@{}r r p{0.16\textwidth} p{0.50\textwidth} c@{}}
\toprule
PU & Printed & Role & Disposition & Note\\
\midrule
\endfirsthead
\toprule
PU & Printed & Role & Disposition & Note\\
\midrule
\endhead
1 & -- & extrinsic wrapper & EXCLUDED; provenance/topology only & 1 \\

2 & unnumbered & title page & ACCEPTED intrinsic text; copy marks excluded & 2 \\

3 & -- & blank witness & BLANK position retained; copy marks excluded & 3 \\

4 & unnumbered & dedication & ACCEPTED & 4 \\

5 & unnumbered & dedication continuation & ACCEPTED & 4 \\

6 & 1 & printed page 1 & ACCEPTED; source-verified mathematical body & -- \\

7 & 2 & printed page 2 & ACCEPTED; source-verified mathematical body & -- \\

8 & 3 & printed page 3 & ACCEPTED; source-verified mathematical body & -- \\

9 & 4 & printed page 4 & ACCEPTED; source-verified mathematical body & -- \\

10 & 5 & printed page 5 & ACCEPTED; source-verified mathematical body & 5 \\

11 & 6 & printed page 6 & ACCEPTED; source-verified mathematical body & -- \\

12 & 7 & printed page 7 & ACCEPTED; source-verified mathematical body & -- \\

13 & 8 & printed page 8 & ACCEPTED; source-verified mathematical body & -- \\

14 & 9 & printed page 9 & ACCEPTED; source-verified mathematical body & -- \\

15 & 10 & printed page 10 & ACCEPTED; source-verified mathematical body & 6 \\

16 & 11 & printed page 11 & ACCEPTED; source-verified mathematical body & 7 \\

17 & 12 & printed page 12 & ACCEPTED; source-verified mathematical body & 8 \\

18 & 13 & printed page 13 & ACCEPTED; source-verified mathematical body & -- \\

19 & 14 & printed page 14 & ACCEPTED; source-verified mathematical body & -- \\

20 & 15 & printed page 15 & ACCEPTED; source-verified mathematical body & -- \\

21 & 16 & printed page 16 & ACCEPTED; source-verified mathematical body & -- \\

22 & 17 & printed page 17 & ACCEPTED; source-verified mathematical body & -- \\

23 & 18 & printed page 18 & ACCEPTED; source-verified mathematical body & 9 \\

24 & 19 & printed page 19 & ACCEPTED; source-verified mathematical body & 10 \\

25 & 20 & printed page 20 & ACCEPTED; source-verified mathematical body & 11 \\

26 & 21 & printed page 21 & ACCEPTED; source-verified mathematical body & -- \\

27 & 22 & printed page 22 & ACCEPTED; source-verified mathematical body & -- \\

28 & 23 & printed page 23 & ACCEPTED; source-verified mathematical body & -- \\

29 & 24 & printed page 24 & ACCEPTED; source-verified mathematical body & -- \\

30 & 25 & printed page 25 & ACCEPTED; source-verified mathematical body & -- \\

31 & 26 & printed page 26 & ACCEPTED; source-verified mathematical body & -- \\

32 & 27 & printed page 27 & ACCEPTED; source-verified mathematical body & -- \\

33 & 28 & printed page 28 & ACCEPTED; source-verified mathematical body & 12 \\

34 & 29 & printed page 29 & ACCEPTED; source-verified mathematical body & -- \\

35 & 30 & printed page 30 & ACCEPTED; source-verified mathematical body & -- \\

36 & 31 & printed page 31 & ACCEPTED; source-verified mathematical body & 13 \\

37 & 32 & printed page 32 & ACCEPTED; source-verified mathematical body & 14 \\

38 & 33 & printed page 33 & ACCEPTED; source-verified mathematical body & -- \\

39 & 34 & printed page 34 & ACCEPTED; source-verified mathematical body & -- \\

40 & 35 & printed page 35 & ACCEPTED; source-verified mathematical body & -- \\

41 & 36 & printed page 36 & ACCEPTED; source-verified mathematical body & 15 \\

42 & 37 & printed page 37 & ACCEPTED; source-verified mathematical body & -- \\

43 & 38 & printed page 38 & ACCEPTED; source-verified mathematical body & 16 \\

44 & 39 & printed page 39 & ACCEPTED; source-verified mathematical body & 17 \\

45 & 40 & printed page 40 & ACCEPTED; source-verified mathematical body & 18 \\

46 & -- & Vita & ACCEPTED intrinsic biography; copy footer excluded & 19 \\

47 & -- & Theses & ACCEPTED intrinsic terminal theses; copy footer excluded & 20 \\

\end{longtable}

\normalsize

\section*{Numbered notes}
\begin{enumerate}
\item[1.] \textbf{PU001 (provider wrapper).} The complete Google Books reproduction notice, logo, QR code, URL, and rule are extrinsic copy/digitizer matter. They remain in the authority evidence and are excluded from both scholarly editions.

\item[2.] \textbf{PU002 (title).} The handwritten accession or shelfmark and the ``Digitized by Google'' footer are copy-specific and excluded. The dissertation title, defence information, opponents, place, and printer are intrinsic and retained. The title rules and central ornament are conservatively reconstructed as editable typographic objects.

\item[3.] \textbf{PU003 (blank witness).} No authorial text is visible. The blank physical position is retained as a blank page; speckle, edge marks, and the digitizer footer are not reproduced.

\item[4.] \textbf{PU004--PU005 (dedication).} The Latin abbreviation \emph{D. D. D.} is retained. The English layer expands its conventional force as ``he gives, dedicates, and devotes'' rather than treating the expansion as Latin source text.

\item[5.] \textbf{PU010 / printed p.5.} The first elliptic-coordinate variational display visibly reads $1-e^2\cos l$, without a superscript on $\cos$. Since nearby formulas use $1-e^2\cos^2 l$, the missing exponent is probably a printing error. The diplomatic Latin and faithful English preserve the printed reading; the probable correction is recorded only here.

\item[6.] \textbf{PU015 / printed p.10.} In the inverse series for $x$, the authority clearly reads $-\frac{33}{32}p^2$. The inherited $p^3$ reading was rejected as a transcription error, and the source reading is restored in both editions.

\item[7.] \textbf{PU016 / printed p.11.} The first-variation symbol carries subscript $1$ and the second variation carries subscript $2$. The inherited omission of the first subscript was rejected. The printed denominator exponent is retained in its mixed form $1\frac12$ rather than silently regularized to $\frac32$.

\item[8.] \textbf{PU017 / printed p.12.} The transformed second-variation denominator again prints the mixed exponent $1\frac12$. The source apostrophe in \emph{integralis'} is retained diplomatically.

\item[9.] \textbf{PU023 / printed p.18.} The dense $U$-reduction was checked directly against the 400-dpi authority replay. In particular, after $u_1=0$ the derivative brace is ordered
\[
 \frac{u_2}{v_2}\left(\frac{\sin\varphi}{v}\right)'_2
 -\frac{\sin\varphi_2}{v_2}\left(\frac{u}{v}\right)'_2.
\]
The accepted editable formula preserves that ordering.

\item[10.] \textbf{PU024 / printed p.19.} The theorem visibly prints $\psi<\psi'$ in its first condition but $\psi>\psi$ in the second. The latter is almost certainly a dropped prime, but it is preserved in the diplomatic and English layers. The conjectural $\psi>\psi'$ correction is confined to this note.

\item[11.] \textbf{PU025 / printed p.20.} The source uses an uppercase Greek $\Pi$ in the interval ``between $0$ and $2\Pi$,'' while the surrounding formulas predominantly use $\pi$. The uppercase form is preserved in the diplomatic layer.

\item[12.] \textbf{PU033 / printed p.28.} Fresh 400-dpi replay establishes the first $E$-series line as $8q+16q^2+32q^3\ldots$. The inherited omission of the exponent in $16q^2$ was rejected in both editions.

\item[13.] \textbf{PU036 / printed p.31.} The extension formula visibly has the coefficient $4q/(1-2q)^3$. A later inherited variant with square rather than cube was rejected.

\item[14.] \textbf{PU037 / printed p.32.} In $\cos A\,\tg o=\tg x$ and in the immediately following right-hand $o+4iq^2\cdots$, the character is a small Latin $o$, not Greek $\sigma$ and not the numeral zero. Both visible occurrences are retained in each edition.

\item[15.] \textbf{PU041 / printed p.36.} The $3\times3$ cosine determinant and the three subsequent determinant reductions are reconstructed as editable mathematics. No diagram or image fallback is used.

\item[16.] \textbf{PU043 / printed p.38.} The text introduces the right triangle $BCD$ and the relation $BD/CD=\tg BCD$ verbally and algebraically; the authority page contains no drawn triangle. No figure is missing from either edition.

\item[17.] \textbf{PU044--PU045 / printed pp.39--40.} The source prints ``e formula:'' once and then uses ditto marks in the left column. The Latin preserves those source-visible marks typographically. The English translates the first phrase as ``from the formula:'' and preserves the subsequent ditto convention.

\item[18.] \textbf{PU045 / printed p.40.} The numerical example visibly prints $x_1=149^\circ38'429{,}65''$. The anomalous seconds value is preserved diplomatically and in the faithful English layer; no normalization to minutes and seconds is made.

\item[19.] \textbf{PU046 (Vita).} Direct replay restores the printed forms \emph{terras exsternas} and the name \emph{Frankenhein}. The English translates the former as ``foreign lands'' and retains the latter name. The Google digitizer footer is copy-specific and excluded; the terminal ornament is reconstructed editably.

\item[20.] \textbf{PU047 (Theses).} All four theses, the title's terminal dot, and the closing ornament are intrinsic and retained. Speckle, edge marks, and the ``Digitized by Google'' footer are copy-specific and excluded.
\end{enumerate}
\end{document}
